\documentclass[11pt, a4paper]{article}
\usepackage[utf8]{inputenc}
\usepackage[italian]{babel}
\usepackage{geometry}
\geometry{margin=1in}
\usepackage{booktabs}
\usepackage{hyperref}
\usepackage{array}
\usepackage{xcolor}

\hypersetup{
    colorlinks=true,
    linkcolor=blue,
    filecolor=magenta,      
    urlcolor=blue,
    pdftitle={Guida Introduttiva a Git},
}

\title{\textbf{GUIDA INTRODUTTIVA A GIT}}
\author{\textbf{Stefano Peris} \\ 
\small Email: \href{mailto:stefanoperis.github@gmail.com}{stefanoperis.github@gmail.com} \\
\small GitHub: \href{https://github.com/TechnoKnight77}{github.com/TechnoKnight77}}
\date{\today}

\begin{document}

\maketitle

\section{Introduzione}
Questa guida rapida raccoglie i comandi essenziali per l'utilizzo di Git, organizzati per logica di utilizzo e pronti per essere consultati durante le sessioni di sviluppo.

\section{Configurazione Iniziale}
Prima di iniziare a tracciare i file, è fondamentale configurare l'identità dell'utente globale per associare correttamente i commit.

\begin{center}
\small
\begin{tabular}{ll}
\toprule
\textbf{Comando} & \textbf{Descrizione} \\
\midrule
\texttt{git config --global user.name "Stefano Peris"} & Imposta il nome utente globale. \\
\texttt{git config --global user.email "stefanoperis.github@gmail.com"} & Imposta l'email globale. \\
\texttt{git config --list} & Mostra tutte le configurazioni attive. \\
\bottomrule
\end{tabular}
\end{center}

\section{Comandi Principali del Workflow}
Il ciclo di vita base del codice all'interno di un repository Git locale e remoto.

\begin{center}
\small
\begin{tabular}{lp{9cm}}
\toprule
\textbf{Comando} & \textbf{Descrizione} \\
\midrule
\texttt{git init} & Inizializza un nuovo repository Git locale nella cartella corrente. \\
\texttt{git clone <url>} & Clona un repository remoto esistente sul proprio computer. \\
\texttt{git status} & Mostra lo stato dei file (modificati, tracciati, in staging). \\
\texttt{git add <file>} & Aggiunge uno specifico file all'area di staging. \\
\texttt{git add .} & Aggiunge tutti i file modificati e nuovi all'area di staging. \\
\texttt{git commit -m "messaggio"} & Registra le modifiche presenti nello staging creandone uno snapshot. \\
\texttt{git push origin <branch>} & Invia i commit locali al repository remoto sul branch specificato. \\
\texttt{git pull} & Recupera le ultime modifiche dal repository remoto e le unisce nel codice locale. \\
\bottomrule
\end{tabular}
\end{center}

\section{Gestione dei Branch (Rami)}
Lavorare in parallelo isolando le funzionalità o i fix di bug.

\begin{center}
\small
\begin{tabular}{lp{9cm}}
\toprule
\textbf{Comando} & \textbf{Descrizione} \\
\midrule
\texttt{git branch} & Elenca tutti i branch locali esistenti. \\
\texttt{git branch <nome-branch>} & Crea un nuovo branch con il nome specificato. \\
\texttt{git checkout <nome-branch>} & Passa al branch specificato lavorando sul suo storico. \\
\texttt{git checkout -b <nome-branch>} & Scorciatoia per creare un nuovo branch e passarvi immediatamente. \\
\texttt{git merge <nome-branch>} & Unisce la cronologia del branch specificato dentro quello attuale. \\
\texttt{git branch -d <nome-branch>} & Elimina in sicurezza il branch specificato (solo se già unito). \\
\bottomrule
\end{tabular}
\end{center}

\end{document}
